%----------------------------------------------------------------------------
\chapter{Introduction}\label{ch:intro}
%----------------------------------------------------------------------------
\guidetodo{Introduction}{app:chapters:intro}

\paragraph{Context.}\label{par:context}
\guidetodo{Context}{app:chapters:context}

\paragraph{Problem.}\label{par:problem}
\guidetodo{Problem}{app:chapters:problem}
\begin{runningexample}
TODO Two or three sentences: the actors, what they want to do, and what goes wrong today. Every chapter returns to this scenario, so choose one small enough to follow and rich enough to show every requirement.
\end{runningexample}

\paragraph{Requirements.}
\guidetodo{Requirements}{app:chapters:requirements}
\begin{requirements}
  \item\label{req:first} \textbf{TODO Short name.} A property the running example needs, stated so it can be checked.
\end{requirements}

\paragraph{Gap.}\label{par:gap}
\guidetodo{Gap}{app:chapters:introgap}

\begin{figure}[!ht]
  \centering
  \missingfigure[figheight=2cm]{Leading figure: \cref{fig:guide:overview}}
  \caption{TODO leading figure.}
  \label{fig:overview}
\end{figure}

\paragraph{Contributions.}
\guidetodo{Contributions}{app:chapters:contributions}
TODO The contributions, which \cref{fig:overview} shows in one picture.
\begin{contributions}
  \item\label{c:method} \textbf{TODO Short name.} A theoretical contribution: a model, formalisation, algorithm or protocol, with what it guarantees (\cref{ch:approach}).
\end{contributions}

\paragraph{Structure.}
\guidetodo{Structure}{app:chapters:structure}
\Cref{ch:background} introduces the background, \cref{ch:approach} presents the approach, \cref{ch:implementation} the implementation, and \cref{ch:evaluation} the evaluation. \Cref{ch:related} compares the results with related work, and \cref{ch:conclusion} concludes with future work.
